% titel: Blatt 6, A4: Wann gilt die Oktettregel nicht?

\subsection*{Frage}
Für welche Verbindungen $\mathrm{CH_4,\ PCl_5,\ BF_3}$, Wasser, $\mathrm{C_5H_12}$
(Pentan) gilt die Oktettregel nicht?

\subsection*{Lösung}
\begin{itemize}
  \item \textbf{$\mathrm{CH_4}$:} gilt. C und jedes H erfüllen ihr Oktett (bzw.
        Duett).
  \item \textbf{$\mathrm{PCl_5}$:} gilt nicht. P (Periode~3) hält 10
        Valenzelektronen (5 Bindungen) --- Oktetterweiterung.
  \item \textbf{$\mathrm{BF_3}$:} gilt nicht. B bleibt beim Elektronenmangel-
        Oktett ($6$ Valenzelektronen, Elektronendefizit).
  \item \textbf{Wasser ($\mathrm{H_2O}$):} gilt. O hat ein Oktett (2
        Bindungen + 2 freie Paare), H ein Duett.
  \item \textbf{$\mathrm{C_5H_{12}}$ (Pentan):} gilt. Alle Atome erfüllen das
        Oktett (kein Elektronenmangel, keine Erweiterung nötig).
\end{itemize}

\begin{ergebnis}
Die Oktettregel gilt \textbf{nicht} für $\mathrm{PCl_5}$ (Erweiterung) und
$\mathrm{BF_3}$ (Elektronenmangel); sie gilt für $\mathrm{CH_4}$, Wasser und
Pentan.
\end{ergebnis}

\begin{kontrolle}
$\mathrm{PCl_5}$ und $\mathrm{BF_3}$ sind beide stabile Stoffe, verletzen aber
die Regel --- der Nachweis ist die Valenzelektronenzahl am Zentralatom (10 bzw.
6). Die drei übrigen Verbindungen sind Oktett-konform, weil jedes Atom (bzw. H)
genau 8 (bzw. 2) Elektronen trägt.
\end{kontrolle}

\begin{merksatz}
Oktett gilt nicht bei: Oktetterweiterung (Perioden $\ge 3$), Elektronenmangel
(Be, B), Radikalen. $\mathrm{CH_4}$, Wasser, Alkane/Kohlenwasserstoffe sind
typisch Oktett-konform.
\end{merksatz}
